\documentclass[../../main.tex]{subfiles}
\begin{document}

{\bf [解]}

$C$の高さを $h = 1600 + x\ \text{[m]}$ とおく．
$B$の高さを基準とすると$A$の高さは
\begin{align}
  y = 1600 - 1210 = 390 \label{eq:1}
\end{align}
である．
$A$, $C$から標高$0m$におろした垂線と，$B$の高さの平面の交点を$A',C'$と置く．
また，$A$から$CC'$におろした垂線の足を$C''$とする．
題意の状況は下図のようになる．

\begin{figure}[htb]
  \centering
  \begin{tikzpicture}[
      scale=1.1,
      >=stealth,
      every node/.style={font=\small},
      x={(0.95cm, -0.22cm)},  % x軸（水平右下方向のプロジェクション）
      y={(0.55cm, 0.42cm)},   % y軸（奥行き方向のプロジェクション）
      z={(0cm, 1.25cm)}       % z軸（鉛直上向き）
    ]

    % --- 基準平面（B の高さ: z = 0）上の頂点 ---
    % B を原点付近に配置
    \coordinate (B)  at (0, 0, 0);
    \coordinate (Ap) at (-3.2, 0.6, 0);   % A'
    \coordinate (Cp) at (2.4, 3.6, 0);    % C'

    % --- 鉛直方向の高さを持つ頂点 ---
    % A の高さ y = 390 -> z = 1.3
    % C'' (C' の真上、A と同じ高さ) -> z = 1.3
    % C の高さ 390 + x -> z = 2.4
    \coordinate (A)  at (-3.2, 0.6, 1.3);
    \coordinate (Cm) at (2.4, 3.6, 1.3);  % C''
    \coordinate (C)  at (2.4, 3.6, 2.4);

    % --- 1. 基準平面 (z = 0) の描画（薄い色で台座を表現） ---
    \fill[gray!10, opacity=0.7] (Ap) -- (B) -- (Cp) -- cycle;
    \draw[thick, gray!80!black] (Ap) -- (B) -- (Cp) -- cycle;

    % --- 2. 鉛直垂線（A-A', C-C'） ---
    \draw[thick, blue!70!black] (Ap) -- (A);
    \draw[thick, blue!70!black] (Cp) -- (Cm);
    \draw[very thick, blue!80!black] (Cm) -- (C);

    % --- 3. 水平補助線 A -- C'' ---
    \draw[thick, dashed,red!70!black] (A) -- (Cm);

    % --- 4. 視線・傾斜線（実線） ---
    \draw[very thick] (A) -- (B);
    \draw[very thick] (B) -- (C);
    \draw[very thick, magenta!90!black] (A) -- (C);

    % --- 5. 直角記号 ---
    % A' における垂直 (BA' \perp AA')
    \draw[gray!70] ($ (Ap) + (0.2, 0, 0) $) -- ($ (Ap) + (0.2, 0, 0.2) $) -- ($ (Ap) + (0, 0, 0.2) $);
    % C' における垂直 (BC' \perp CC')
    \draw[gray!70] ($ (Cp) + (-0.2, 0, 0) $) -- ($ (Cp) + (-0.2, 0, 0.2) $) -- ($ (Cp) + (0, 0, 0.2) $);
    % C'' における垂直 (AC'' \perp C''C)
    \draw[gray!70] ($ (Cm) + (-0.2, 0, 0) $) -- ($ (Cm) + (-0.2, 0, 0.2) $) -- ($ (Cm) + (0, 0, 0.2) $);

    % --- 6. 仰角の円弧表示 ---
    % \angle ABA' = 30°
    \draw[thick, green!50!black] ($ (B)!0.28!(Ap) $) to[bend left=12] ($ (B)!0.28!(A) $);
    \node[green!50!black, left=1pt, font=\footnotesize] at ($ (B)!0.35!(A) + (-0.1, 0, -0.05) $) {$30^\circ$};

    % \angle CBC' = 30°
    \draw[thick, green!50!black] ($ (B)!0.28!(Cp) $) to[bend left=-15] ($ (B)!0.28!(C) $);
    \node[green!50!black, right=1pt, font=\footnotesize] at ($ (B)!0.32!(C) + (0.05, 0, 0) $) {$30^\circ$};

    % \angle CAC'' = 15° (仰角)
    \draw[thick, orange!90!black] ($ (A)!0.28!(Cm) $) to[bend right=15] ($ (A)!0.28!(C) $);
    \node[orange!90!black, above right=1pt, font=\footnotesize] at ($ (A)!0.32!(C) $) {$15^\circ$};

    % --- 7. 各点のプロットとラベル ---
    \fill (B) circle (1.5pt) node[below=2pt] {$B$};
    \fill (Ap) circle (1.3pt) node[below left] {$A'$};
    \fill (Cp) circle (1.3pt) node[below right] {$C'$};
    \fill (A) circle (1.5pt) node[above left] {$A$};
    \fill (Cm) circle (1.3pt) node[right=2pt] {$C''$};
    \fill (C) circle (1.5pt) node[above=2pt] {$C$};

    % --- 8. 高さ寸法のラベル ---
    \node[left, blue!80!black, font=\footnotesize] at ($ (Ap)!0.5!(A) $) {$390$};
    \node[right, blue!80!black, font=\footnotesize] at ($ (Cp)!0.5!(Cm) $) {$390$};
    \node[right, blue!80!black, font=\footnotesize] at ($ (Cm)!0.5!(C) $) {$x$};

  \end{tikzpicture}
  \caption{点 $A, B, C$ の立体的な位置関係と仰角}
  \label{fig:3d_abc_heights}
\end{figure}

$A$から$C$への仰角は$15^\circ$であり，$\tan 15^\circ$ は加法定理より
\begin{align}
  \tan 15^{\circ}
  = \tan \left(45^{\circ}-30^{\circ}\right)
  = \dfrac{\tan 45^{\circ} - \tan 30^{\circ}}{1+\tan 45^{\circ} \tan 30^{\circ}}
  = \dfrac{1-1/\sqrt{3}}{1+1/\sqrt{3}}
  = 2-\sqrt{3} \label{eq:2}
\end{align}
である．
従って各辺の長さは
\begin{align}
  |BC'|  &= \dfrac{|CC'|}{\tan \angle CBC'} = \dfrac{|CC'|}{\tan 30^{\circ}}  = \sqrt{3}(x+y) \\
  |BA'|  &= \dfrac{|AA'|}{\tan \angle ABA'} = \dfrac{|AA'|}{\tan 30^{\circ}}  = \sqrt{3}y \\
  |A'C'| &= |AC''| = \dfrac{|CC''|}{\tan \angle CAC''} = \dfrac{|CC''|}{\tan 15^{\circ}} = \dfrac{x}{2-\sqrt{3}} = \left(2+\sqrt{3}\right)x
\end{align}
である．よって$\triangle A'BC'$は下図のようになる．

\begin{figure}[htb]
  \centering
  \begin{tikzpicture}[scale=1.2, >=stealth, every node/.style={font=\small}]

    % --- 1. 座標定義 ---
    \coordinate (Cp) at (1.5, 4.0);
    \coordinate (Ap) at ($ (Cp) + (240:3.2) $);
    \coordinate (B)  at ($ (Cp) + (300:3.8) $);

    % --- 2. 三角形 A'BC' の描画 ---
    \draw[very thick] (Ap) -- (B) -- (Cp) -- cycle;

    % --- 3. 鉛直補助線（破線） ---
    \coordinate (Ap_top) at ($ (Ap) + (0, 1.4) $);
    \draw[thick, dashed, gray!80!black] (Ap) -- (Ap_top);

    \coordinate (B_top) at ($ (B) + (0, 1.6) $);
    \draw[thick, dashed, gray!80!black] (B) -- (B_top);

    % --- 4. 角度の円弧とラベル（構文修正済み） ---
    % (1) A' における角度 10°
    \draw[thick, red!80!black] (Ap) +(90:0.8) arc[start angle=90, end angle=60, radius=0.8];
    \node[red!80!black, font=\footnotesize] at ($(Ap) + (75:1.05)$) {$10^\circ$};

    % (2) B における角度 20°
    \draw[thick, red!80!black] (B) +(90:0.9) arc[start angle=90, end angle=120, radius=0.9];
    \node[red!80!black, font=\footnotesize] at ($(B) + (105:1.15)$) {$20^\circ$};

    % (3) C' における頂角 30°
    \draw[thick, blue!80!black] (Cp) +(240:0.7) arc[start angle=240, end angle=300, radius=0.7];
    \node[blue!80!black, font=\footnotesize] at ($(Cp) + (270:0.95)$) {$30^\circ$};

    % --- 5. 頂点プロットとラベル ---
    \fill (Cp) circle (1.3pt) node[above=2pt] {$C'$};
    \fill (Ap) circle (1.3pt) node[below left=1pt] {$A'$};
    \fill (B)  circle (1.3pt) node[below right=1pt] {$B$};

    % --- 6. 辺の長さラベル ---
    \node[above left, font=\footnotesize] at ($ (Ap)!0.5!(Cp) + (-0.15, 0.1) $) {$(2+\sqrt{3})x$};
    \node[above right, font=\footnotesize] at ($ (B)!0.5!(Cp) + (0.1, 0) $) {$\sqrt{3}(x+y)$};
    \node[below=3pt, font=\footnotesize] at ($ (Ap)!0.5!(B) $) {$\sqrt{3}y$};

  \end{tikzpicture}
  \caption{基準平面上の $\triangle A'BC'$ と各点の方位角・辺の長さ}
  \label{fig:triangle_ap_b_cp_bearings}
\end{figure}

また，題意より$\angle A'C'B=10^{\circ}+20^{\circ}=30^{\circ}$である．
だから $\triangle A'BC'$ に余弦定理を用いて
\begin{align}
  & |A'B|^2 = |A'C'|^2 + |BC'|^2 - 2|A'C'||BC'|\cos\angle A'C'B \\
  \therefore
  & \left(\sqrt{3}y\right)^2 = \left(2+\sqrt{3}\right)^2x^2 + 3\left(x+y\right)^2 - 2\sqrt{3}\left(2+\sqrt{3}\right)x(x+y)\dfrac{\sqrt{3}}{2}
\end{align}
である．整理して
\begin{align}
  (\sqrt{3}y)^2 &= \Bigl((2+\sqrt{3})x\Bigr)^2 + 3(x+y)^2 - 2\sqrt{3}(2+\sqrt{3})x(x+y)\dfrac{\sqrt{3}}{2} \\
  3A^2 &= (7+4\sqrt{3})x^2 + 3(x^2 + 2yx + y^2) - (6+3\sqrt{3})(x^2 + xy) \\
  0 &= \Bigl\{(7+4\sqrt{3}) + 3 - (6+3\sqrt{3})\Bigr\}x^2 + \left\{6y - \left(6 + 3\sqrt{3}\right)y\right\}x \\
  0 &= (4+\sqrt{3})x^2 - 3\sqrt{3}yx
\end{align}
を得る．両辺を$x \neq 0$ で割って
\begin{align}
  (4+\sqrt{3})x &= 3\sqrt{3}y
\end{align}
である．従って$x$は
\begin{align}
  x
  &= \dfrac{3\sqrt{3}}{4+\sqrt{3}}y \\
  &= \dfrac{-9+12\sqrt{3}}{13}y \\
  &= \dfrac{-9+12\cdot 1.732}{13} \cdot 390 \\
  &= (-9+20.784) \cdot 30 \\
  &= 353.52
\end{align}
だから，求める高さ$h$は
\begin{align}
  h = 1600 + 353.52 = 1953.52 \fallingdotseq 1954\text{ [m]}
\end{align}
である．

\end{document}
